% from ega1-10-fr.tex
\oldpage[I]{202}

Notons que la propriété d'exactitude de $M^\Delta$, appliqué à la suite
exacte
$0\to\mathfrak{J}\to A\to A/\mathfrak{J}\to0$ montre que le faisceau
d'idéaux de $\mathscr{O}_{\mathfrak{X}}$ désigné ici par
$\mathfrak{J}^\Delta$ coïncide avec celui qui avait été noté de la même
façon en (\hyperref[I.10.3.1-fr]{10.3.1}), en vertu de
(\hyperref[I.10.3.2-fr]{10.3.2}).

\begin{proposition}[10.10.3]
\label{I.10.10.3-fr}
Sous les hypothèses de (\hyperref[I.10.10.1-fr]{10.10.1}),
$\mathscr{O}_{\mathfrak{X}}$ est un faisceau cohérent d'anneaux.
\end{proposition}

Si $f\in A$, on sait que $A_{\{f\}}$ est un anneau adique noethérien
(\hyperref[0.7.6.11-fr]{0, 7.6.11}) et comme la question est locale, on
est ramené (\hyperref[I.10.1.4-fr]{10.1.4}) à prouver que le noyau d'un
homomorphisme
$v:\mathscr{O}_{\mathfrak{X}}^n\to\mathscr{O}_{\mathfrak{X}}$ est un
$\mathscr{O}_{\mathfrak{X}}$-Module de type fini. On a alors
$v=u^\Delta$, où $u$ est un $A$-homomorphisme $A^n\to A$
(\hyperref[I.10.10.2-fr]{10.10.2})~; comme $A$ est noethérien, le noyau
de $u$ est de type fini, autrement dit on a un homomorphisme
$A^m\xrightarrow{w}A^n$ tel que la suite
$A^m\xrightarrow{w}A^n\xrightarrow{u}A$ soit exacte. On en conclut
(\hyperref[I.10.10.2-fr]{10.10.2}) que la suite
$\mathscr{O}_{\mathfrak{X}}^m\xrightarrow{w^\Delta}
\mathscr{O}_{\mathfrak{X}}^n\xrightarrow{v}
\mathscr{O}_{\mathfrak{X}}$ est exacte, ce qui prouve que le noyau de
$v$ est de type fini.

\begin{env}[10.10.4]
\label{I.10.10.4-fr}
Avec les notations précédentes, posons
$A_n=A/\mathfrak{J}^{n+1}$, et soit $X_n$ le schéma affine
$\operatorname{Spec}(A_n)=(\mathfrak{X},
\mathscr{O}_{\mathfrak{X}}/\mathscr{J}^{n+1})$,
$\mathscr{J}=\mathfrak{J}^\Delta$ étant le faisceau d'idéaux de
définition de $\mathscr{O}_{\mathfrak{X}}$ correspondant à l'idéal
$\mathfrak{J}$. Soit $u_{mn}$ le morphisme de préschémas
$X_m\to X_n$ correspondant à l'homomorphisme canonique
$A_n\to A_m$ pour $m\leq n$~; le schéma formel $\mathfrak{X}$ est
limite inductive des $X_n$ pour les $u_{mn}$
(\hyperref[I.10.6.3-fr]{10.6.3}).
\end{env}

\begin{proposition}[10.10.5]
\label{I.10.10.5-fr}
Sous les hypothèses de (\hyperref[I.10.10.1-fr]{10.10.1}), soit
$\mathscr{F}$ un $\mathscr{O}_{\mathfrak{X}}$-Module. Les conditions
suivantes sont équivalentes~:
\begin{enumerate}
  \item[a)] $\mathscr{F}$ est un $\mathscr{O}_{\mathfrak{X}}$-Module
    cohérent.
  \item[b)] $\mathscr{F}$ est isomorphe à la limite projective
    (\hyperref[I.10.6.6-fr]{10.6.6}) d'une suite $(\mathscr{F}_n)$ de
    $\mathscr{O}_{X_n}$-Modules cohérents tels que
    $u_{mn}^*(\mathscr{F}_n)=\mathscr{F}_m$.
  \item[c)] Il existe un $A$-module de type fini $M$ (déterminé à un
    isomorphisme canonique près par
    (\hyperref[I.10.10.2-fr]{10.10.2, (i)})) tel que $\mathscr{F}$ soit
    isomorphe à $M^\Delta$.
\end{enumerate}
\end{proposition}

Montrons d'abord que b) implique c). On a
$\mathscr{F}_n=\widetilde{M_n}$, où $M_n$ est un $A_n$-module de type
fini, et l'hypothèse entraîne que
$M_m=M_n\otimes_{A_n}A_m$ pour $m\leq n$
(\hyperref[I.1.6.5-fr]{1.6.5})~; les $M_n$ forment donc un système
projectif pour les di-homomorphismes canoniques $M_n\to M_m$
($m\leq n$), et il résulte aussitôt de la définition des $A_n$ que ce
système projectif vérifie les conditions de
(\hyperref[0.7.2.9-fr]{0, 7.2.9})~; sa limite projective $M$ est par
suite un $A$-module de type fini tel que
$M_n=M\otimes_A A_n$ pour tout $n$. On en déduit que $\mathscr{F}_n$
est induit sur $X_n$ par
$\widetilde M\otimes_{\mathscr{O}_X}
(\mathscr{O}_X/\widetilde{\mathfrak{J}}^{n+1})$, donc
$\mathscr{F}=M^\Delta$ par définition
(\hyperref[I.10.8.4-fr]{10.8.4}).

Inversement, c) entraîne b)~; en effet, si $u_n$ est le morphisme
d'immersion $X_n\to X$,
$u_n^*(\widetilde M)=(M\otimes_A A_n)^{\sim}$ est induit sur $X_n$ par
$\widetilde M\otimes_{\mathscr{O}_X}
(\mathscr{O}_X/\widetilde{\mathfrak{J}}^{n+1})$, et
$M^\Delta=\varprojlim u_n^*(\widetilde M)$ par définition
(\hyperref[I.10.8.4-fr]{10.8.4})~; comme
$u_m=u_n\circ u_{mn}$ pour $m\leq n$, les
$\mathscr{F}_n=u_n^*(\widetilde M)$ vérifient les conditions de b),
d'où notre assertion.

Montrons maintenant que c) implique a)~: en effet, on a par définition
$\mathscr{O}_{\mathfrak{X}}=A^\Delta$~; $M$ étant conoyau d'un
homomorphisme $A^m\to A^n$, il résulte de
(\hyperref[I.10.10.2-fr]{10.10.2}) que $M^\Delta$ est conoyau d'un
homomorphisme
$\mathscr{O}_{\mathfrak{X}}^m\to\mathscr{O}_{\mathfrak{X}}^n$, et
comme le faisceau d'anneaux $\mathscr{O}_{\mathfrak{X}}$ est cohérent
(\hyperref[I.10.10.3-fr]{10.10.3}), il en est de même de $M^\Delta$
(\hyperref[0.5.3.4-fr]{0, 5.3.4}).
